\section{Redundancy and Primary Copy Identification}

Cloud storage is, by default, redundant: providers routinely replicate objects across multiple
physical locations (and often across regions) for durability and disaster-recovery purposes, with
no action required from the customer. From a forensic standpoint this redundancy is a
double-edged sword: it means the data is unlikely to be permanently lost, but it also means that
``the same file'' may legitimately exist in more than one place at once, at slightly different
points in a replication cycle --- and only one of those copies was actually the one interacted
with at the time of the incident.

\subsection{Inspecting the Version Export}

\begin{lstlisting}
jq '.versions[] | {versionId, region, lastModified, etag, replicationStatus}' s3_object_versions.json
\end{lstlisting}

\begin{questionbox}
Three entries are listed for \texttt{contract\_2026\_0091.pdf}. Which one is the primary,
authoritative copy at the time of the incident, and which is a replica? Justify your answer
using at least two independent fields from the export --- do not rely on \texttt{lastModified}
alone.
\end{questionbox}

\begin{answerbox}
\vspace{3cm}
\end{answerbox}

\subsection{The Replication Trap}

\begin{questionbox}
A colleague on your team argues that the \texttt{us-east-1} entry must be the ``real''
exfiltrated version, since it has the most recent timestamp. Explain why this reasoning is
flawed, and what evidence in the export directly refutes it.
\end{questionbox}

\begin{answerbox}
\vspace{3cm}
\end{answerbox}

\subsection{Chain of Custody Documentation}

\begin{questionbox}
Draft the chain-of-custody entry (in the same format used in Laboratory 2, Section 6.2) for
the primary copy identified above. What field(s), specific to a cloud/versioned-object context,
do you need to record here that had no equivalent in the Laboratory 2 disk-image chain of
custody?
\end{questionbox}

\begin{answerbox}
\vspace{3.5cm}
\end{answerbox}
